The goal of this section is to describe how the microcosm principle allows one to define general enriched slice categories, especially under the consideration of model structures. Some special cases of the following results have already appeared in the literature for slice-over categories. We provide a general characterization encompassing both slice-over and slice-under categories.

\subsection{Slice monoidal categories and module} \label{sub:slicecaten}

The definition of a module category is again within the scope of the microcosm principle. Consequentially, we start with the definition of a module object inside a module category.

\begin{definition}
	Let $(\cC,\triangleleft_{\cC},a_{\cC},\lambda_{\cC})$ be a left $\cV$-module category over some monoidal category $\cV$. Let $(M,\mu,e) \in \cV$ be a monoid object. We say that an object $Z \in \cC$ is a left module object over $M$ if it has following data:
	\begin{enumerate}
		\item a left action morphism $\triangleleft_{Z} \colon M \triangleleft_{\cC} Z \rightarrow Z$,
		\item the associativity and unitality axioms hold
		% https://q.uiver.app/#q=WzAsOCxbMCwyLCIoTSBcXG90aW1lcyBNKSBcXHRyaWFuZ2xlbGVmdF97XFxjQ30gWiJdLFswLDAsIk0gXFx0cmlhbmdsZWxlZnRfe1xcY0N9IChNIFxcdHJpYW5nbGVsZWZ0X3tcXGNDfSBaKSJdLFswLDQsIk0gXFx0cmlhbmdsZWxlZnRfe1xcY0N9IFoiXSxbMiwwLCJNIFxcdHJpYW5nbGVsZWZ0X3tcXGNDfSBaIl0sWzIsNCwiWiJdLFszLDAsIkkgXFx0cmlhbmdsZWxlZnRfe1xcY0N9IFoiXSxbNSwwLCJNIFxcdHJpYW5nbGVsZWZ0X3tcXGNDfSBaIl0sWzUsMiwiWiJdLFswLDEsImFfe1xcY0N9Il0sWzAsMiwiXFxtdSBcXHRyaWFuZ2xlbGVmdCAxIiwyXSxbMSwzLCIxIFxcdHJpYW5nbGVsZWZ0X3tcXGNDfSAoXFx0cmlhbmdsZWxlZnRfe1p9KSJdLFszLDQsIlxcdHJpYW5nbGVsZWZ0X3tafSJdLFsyLDQsIlxcdHJpYW5nbGVsZWZ0X3tafSIsMl0sWzUsNiwiZSBcXHRyaWFuZ2xlbGVmdF97XFxjQ30gWiJdLFs2LDcsIlxcdHJpYW5nbGVsZWZ0X3tafSJdLFs1LDcsIlxcbGFtYmRhX3tcXGNDfSIsMl1d
		\[\begin{tikzcd}[ampersand replacement=\&]
			{M \triangleleft_{\cC} (M \triangleleft_{\cC} Z)} \&\& {M \triangleleft_{\cC} Z} \& {I \triangleleft_{\cC} Z} \&\& {M \triangleleft_{\cC} Z} \\
			\\
			{(M \otimes M) \triangleleft_{\cC} Z} \&\&\&\&\& Z \\
			\\
			{M \triangleleft_{\cC} Z} \&\& Z
			\arrow["{1 \triangleleft_{\cC} (\triangleleft_{Z})}", from=1-1, to=1-3]
			\arrow["{\triangleleft_{Z}}", from=1-3, to=5-3]
			\arrow["{e \triangleleft_{\cC} Z}", from=1-4, to=1-6]
			\arrow["{\lambda_{\cC}}"', from=1-4, to=3-6]
			\arrow["{\triangleleft_{Z}}", from=1-6, to=3-6]
			\arrow["{a_{\cC}}", from=3-1, to=1-1]
			\arrow["{\mu \triangleleft 1}"', from=3-1, to=5-1]
			\arrow["{\triangleleft_{Z}}"', from=5-1, to=5-3]
		\end{tikzcd}\cdpunct[8pt]{.}\]
	\end{enumerate}
	
\end{definition}

\begin{definition}
	Let $(\cC,\triangleright_{\cC},a_{\cC},\rho_{\cC})$ be a right $\cV$-module category over some monoidal category $\cV$. Let $(M,\mu,e) \in \cV$ be a monoid object. We say that an object $Z \in \cC$ is a right module object over $M$ if it has following data:
	\begin{enumerate}
		\item a right action morphism $\triangleright_{Z} \colon Z \triangleright_{\cC} M \rightarrow Z$,
		\item the associativity and unitality axioms hold
		% https://q.uiver.app/#q=WzAsOCxbMCwyLCJaIFxcdHJpYW5nbGVyaWdodF97XFxjQ30oTSBcXG90aW1lcyBNKSAiXSxbMCwwLCIoWiBcXHRyaWFuZ2xlcmlnaHRfe1xcY0N9IE0pIFxcdHJpYW5nbGVyaWdodF97XFxjQ30gTSJdLFswLDQsIlogXFx0cmlhbmdsZXJpZ2h0X3tcXGNDfSBNIl0sWzIsMCwiWiBcXHRyaWFuZ2xlbGVmdF97XFxjQ30gTSJdLFsyLDQsIloiXSxbMywwLCJaIFxcdHJpYW5nbGVyaWdodF97XFxjQ30gSSJdLFs1LDAsIlogXFx0cmlhbmdsZXJpZ2h0X3tcXGNDfSBNIl0sWzUsMiwiWiJdLFsxLDAsImFfe1xcY0N9IiwyXSxbMCwyLCIxIFxcdHJpYW5nbGVyaWdodF97XFxjQ30gXFxtdSIsMl0sWzEsMywiKFxcdHJpYW5nbGVyaWdodF97Wn0pIFxcdHJpYW5nbGVyaWdodF97XFxjQ30gMSJdLFszLDQsIlxcdHJpYW5nbGVyaWdodF97Wn0iXSxbMiw0LCJcXHRyaWFuZ2xlcmlnaHRfe1p9IiwyXSxbNSw2LCJaIFxcdHJpYW5nbGVyaWdodF97XFxjQ30gZSJdLFs2LDcsIlxcdHJpYW5nbGVyaWdodF97Wn0iXSxbNSw3LCJcXHJob197XFxjQ30iLDJdXQ==
		\[\begin{tikzcd}[ampersand replacement=\&,cramped]
			{(Z \triangleright_{\cC} M) \triangleright_{\cC} M} \&\& {Z \triangleleft_{\cC} M} \& {Z \triangleright_{\cC} I} \&\& {Z \triangleright_{\cC} M} \\
			\\
			{Z \triangleright_{\cC}(M \otimes M) } \&\&\&\&\& Z \\
			\\
			{Z \triangleright_{\cC} M} \&\& Z
			\arrow["{(\triangleright_{Z}) \triangleright_{\cC} 1}", from=1-1, to=1-3]
			\arrow["{a_{\cC}}"', from=1-1, to=3-1]
			\arrow["{\triangleright_{Z}}", from=1-3, to=5-3]
			\arrow["{Z \triangleright_{\cC} e}", from=1-4, to=1-6]
			\arrow["{\rho_{\cC}}"', from=1-4, to=3-6]
			\arrow["{\triangleright_{Z}}", from=1-6, to=3-6]
			\arrow["{1 \triangleright_{\cC} \mu}"', from=3-1, to=5-1]
			\arrow["{\triangleright_{Z}}"', from=5-1, to=5-3]
		\end{tikzcd}\cdpunct[8pt]{.}\]
	\end{enumerate}
\end{definition}

A similar definition can be stated for bimodule object in a bimodule category. One only needs to add a balancing axiom for the bimodule object relative to the balancing $2$-cell of bimodule category. 

As expected, these definitions also provides the slice constructions. Let $(\cV,\otimes,I,a_{\cV},\lambda_{\cV},\rho_{\cV})$ be a monoidal category in what follows.

\begin{lemma} \label{lem:slicingmonandmod}
	Let $(\cC,\triangleleft_{\cC},a_{\cC},\lambda_{\cC})$ be a left $\cV$-module. Let $(M, \mu, e)$ be a monoid object in $\cV$, and let $(Z, \triangleleft_{Z})$ be a left $M$-module. 
	\begin{enumerate}
		\item The slice category $\cV/M$ admits a monoidal structure denoted by $(\cV/M, \otimes_{/M}, I_{/M})$. If $\cV$ is braided or symmetric and $M$ is commutative then $\cV/M$ is also braided or symmetric. 
		\item There is a canonical left $\cV/M$-module structure on $\cC/Z$ denoted by $(\cC/Z,\triangleleft_{/})$.
	\end{enumerate}
\end{lemma}

\begin{proof}
	For the first statement, the sliced tensoring operation is defined by 
	\[f \otimes_{/M} g \coloneqq A \otimes B \xrightarrow{f \otimes g} M \otimes M \xrightarrow{\mu} M\] 
	and the unit by \[I_{/M} \coloneqq I \xrightarrow{\eta} M.\] 
	The naturality of the associator and the unitors of $\cV$ together with the monoid axioms of $M$ allows us to extend these given associator and unitors to those of $\cV/M$ directly.
	
	For the second statement, the sliced left action operation is defined exactly the same. Let $\alpha \colon X \rightarrow M \in \cV/M$ and $f \colon A \rightarrow Z \in \cC/Z$. The sliced left action is defined by \[\alpha \triangleleft_{/} f \colon X \triangleleft A \xrightarrow{\alpha \triangleleft f} M \triangleleft Z \xrightarrow{\triangleleft_{Z}} Z.\] 
	Thereafter, the extension of coherence $2$-cells follow again by the naturality and module axioms.	
\end{proof}

This lemma can be dualized in two ways: one by using the $\op_1$ duality and the other by $r$ duality. The reversal duality, however, provides nothing new.

For the $\op_1$ duality, we consider a monoid object $N$ in the monoidal category $\cV^{\op}$ and a module object $Z$ in $\cC^{\op}$. Applying the previous lemma in this situation and transporting structure back using $\op_1$ duality, we obtain the following.

\begin{lemma}
	Let $(\cC,\triangleleft_{\cC},a_{\cC},\lambda_{\cC})$ be a left $\cV$-module. Let $(N, \mu, \eta)$ be a comonoid object and $(Z,\lambda)$ be a left $N$-comodule. 
	\begin{enumerate}
		\item The coslice category $N/\cV$ admits a monoidal structure denoted by $(N/\cV, \otimes_{N/}, I_{N/})$. If $\cV$ is braided or symmetric and $N$ is commutative then $N/\cV$ is also braided or symmetric. 
		\item There is a canonical left $N/\cV$-module structure on $Z/\cV$ denoted by $(Z/\cC,\triangleleft_{Z/\cC})$.
	\end{enumerate}
\end{lemma}

We now investigate the presence of closures on left $\cV$-module structure. That is, we investigate enrichment through tensoring or powering. Consequentially, we shall now refer to the action functor $\triangleleft \colon \cV \times \cC \rightarrow \cC$ as \emph{tensoring}.

\begin{definition}
	A left $\cV$-module category $(\cC,\triangleleft)$ is called left closed if for every $c \in \cC$, the functor $- \triangleleft c \colon \cV \rightarrow \cV$ has a right adjoint $\ucC(c,-)\colon \cC \rightarrow \cV$. Thus, we have a functor $\ucC(-,-) \colon \cC^{\op} \times \cC \rightarrow \cV$, which is called the \emph{internal hom} of $\cC$.
\end{definition}

\begin{definition}
	A left $\cV$-module category $(\cC,\triangleleft)$ is called right closed if for every $x \in \cV$, the functor $x \triangleleft - \colon \cC \rightarrow \cC$ has a right adjoint $- \pitchfork x \colon \cC \rightarrow \cC$. Thus, we have a functor $- \pitchfork - \colon \cC \times \cV^{\op} \rightarrow \cC$, which is called \emph{powering}. This is usually also denoted by $b^x$ or $\{x,b\}$.
\end{definition}

\begin{definition}
	A left $\cV$-module category $(\cC,\triangleleft)$ is closed if it is both left closed and right closed.
\end{definition}

\begin{lemma}
	A closed left $\cV$-module category $(\cC,\triangleleft)$ is equivalently specified by the data of a two variable left adjunction 
	%% https://q.uiver.app/#q=WzAsOSxbMCwwLCIoLSBcXHRyaWFuZ2xlbGVmdCAtIClee1xcb3B9XFxjb2xvbiJdLFswLDEsIlxcdWNDKC0sLSkgXFxjb2xvbiJdLFswLDIsIi0gXFxwaXRjaGZvcmsgLSBcXGNvbG9uIl0sWzEsMiwiXFxjQyBcXHRpbWVzIFxcY1Zee1xcb3B9Il0sWzEsMSwiXFxjQ157XFxvcH0gXFx0aW1lcyBcXGNDIl0sWzEsMCwiXFxjVl57XFxvcH0gXFx0aW1lcyBcXGNDXntcXG9wfSJdLFsyLDAsIlxcY0Nee1xcb3B9Il0sWzIsMSwiXFxjViJdLFsyLDIsIlxcY0MiXSxbNSw2XSxbNCw3XSxbMyw4XV0=
	\[\begin{tikzcd}[ampersand replacement=\&,cramped]
		{(- \triangleleft - )^{\op}\colon} \& {\cV^{\op} \times \cC^{\op}} \& {\cC^{\op}} \\
		{\ucC(-,-) \colon} \& {\cC^{\op} \times \cC} \& \cV \\
		{- \pitchfork - \colon} \& {\cC \times \cV^{\op}} \& \cC
		\arrow[from=1-2, to=1-3]
		\arrow[from=2-2, to=2-3]
		\arrow[from=3-2, to=3-3]
	\end{tikzcd}\cdpunct[2pt]{.}\]
\end{lemma}
\begin{proof}
	See \cite[Def.~2.1 and Sec.~2.2]{chengmates2014} for the definition and the respective proof. In simpler terms, this is equivalent to the natural isomorphisms $\cC(x \triangleleft a,b) \cong \cV(x,\ucC(a,b)) \cong \cC(a,b \pitchfork x)$.
\end{proof}

\begin{example}
	If the canonical left $\cV$-module $(\cV, \otimes)$ is left closed then it is a left closed monoidal category in the usual sense. The same is true for the right closure of this left module structure.
\end{example}

\begin{lemma}
	There is a canonical enrichment of $\cC$ over $\cV$ using the internal homs $\ucC(-,-)$, which we call $\ucC$. Furthermore, the underlying category of $\ucC$ is isomorphic to $\cC$ itself.
\end{lemma}

\begin{proof}
	One can mimic the proof provided in \cite[Vol.~III, Sec.~3.1]{yau2024bimodcats} by defining $\mathrm{eval}^{c}_d \colon \ucC(c,d) \triangleleft c \rightarrow d$ through the counit and $\mathrm{coeval}^{d}_x \colon x \rightarrow \ucC(d,x \triangleleft d)$ through the unit. Thus, compositions are given by the adjoint to the composites
	% https://q.uiver.app/#q=WzAsNCxbMCwwLCIoXFxjQyhkLGUpIFxcb3RpbWVzIFxcY0MoYyxkKSkgXFx0cmlhbmdsZWxlZnQgYyJdLFswLDEsIlxcY0MoZCxlKSBcXHRyaWFuZ2xlbGVmdChcXGNDKGMsZCkgXFx0cmlhbmdsZWxlZnQgYykiXSxbMywxLCJcXGNDKGQsZSkgXFx0cmlhbmdsZWxlZnQgZCJdLFszLDAsImUiXSxbMCwxLCJhIl0sWzEsMiwiMSBcXHRyaWFuZ2xlbGVmdCBcXG1hdGhybXtldmFsfV9kXntjfSIsMl0sWzIsMywiXFxtYXRocm17ZXZhbH1fZV57ZH0iLDJdLFswLDNdXQ==
	\[\begin{tikzcd}[ampersand replacement=\&,cramped]
		{(\cC(d,e) \otimes \cC(c,d)) \triangleleft c} \&\&\& e \\
		{\cC(d,e) \triangleleft(\cC(c,d) \triangleleft c)} \&\&\& {\cC(d,e) \triangleleft d}
		\arrow[from=1-1, to=1-4]
		\arrow["a", from=1-1, to=2-1]
		\arrow["{1 \triangleleft \mathrm{eval}_d^{c}}"', from=2-1, to=2-4]
		\arrow["{\mathrm{eval}_e^{d}}"', from=2-4, to=1-4]
	\end{tikzcd}\cdpunct[8pt]{,}\]
	and unitors are given by the adjoint to the morphisms $I \triangleleft x \xrightarrow{\lambda_{x}} x$. Thereafter, the coherence follows from the $2$-cells of the left $\cV$-module structure. 
\end{proof}

Let $i\colon I \rightarrow \ucC(c,c)$ be the unitors and let $m\colon \ucC(d,e) \otimes \ucC(c,d) \rightarrow \ucC(c,e)$ be the compositors. In particular, we obtain the following two commuting diagrams
% https://q.uiver.app/#q=WzAsOCxbMCwwLCIoXFxjQyhkLGUpIFxcb3RpbWVzIFxcY0MoYyxkKSkgXFx0cmlhbmdsZWxlZnQgYyJdLFswLDEsIlxcY0MoZCxlKSBcXHRyaWFuZ2xlbGVmdChcXGNDKGMsZCkgXFx0cmlhbmdsZWxlZnQgYykiXSxbMCwyLCJcXGNDKGQsZSkgXFx0cmlhbmdsZWxlZnQgZCJdLFsyLDIsImUiXSxbMiwwLCJcXHVjQyhjLGUpIFxcdHJpYW5nbGVsZWZ0IGMiXSxbMywwLCJJIFxcdHJpYW5nbGVsZWZ0IGMiXSxbNSwyLCJjIl0sWzUsMCwiXFx1Y0MoYyxjKSAgXFx0cmlhbmdsZWxlZnQgYyJdLFswLDEsImEiXSxbMSwyLCIxIFxcdHJpYW5nbGVsZWZ0IFxcbWF0aHJte2V2YWx9X2Ree2N9IiwyXSxbMiwzLCJcXG1hdGhybXtldmFsfV9lXntkfSIsMl0sWzAsNCwibSBcXHRyaWFuZ2xlbGVmdCAxIl0sWzQsMywiXFxtYXRocm17ZXZhbH1fZV57Y30iXSxbNSw2LCJcXGxhbWJkYSIsMl0sWzcsNiwiXFxtYXRocm17ZXZhbH1eY19jIl0sWzUsNywiaSBcXHRyaWFuZ2xlbGVmdCAxIl1d
\[\begin{tikzcd}[ampersand replacement=\&,cramped]
	{(\cC(d,e) \otimes \cC(c,d)) \triangleleft c} \&\& {\ucC(c,e) \triangleleft c} \& {I \triangleleft c} \&\& {\ucC(c,c)  \triangleleft c} \\
	{\cC(d,e) \triangleleft(\cC(c,d) \triangleleft c)} \\
	{\cC(d,e) \triangleleft d} \&\& e \&\&\& c
	\arrow["{m \triangleleft 1}", from=1-1, to=1-3]
	\arrow["a", from=1-1, to=2-1]
	\arrow["{\mathrm{eval}_e^{c}}", from=1-3, to=3-3]
	\arrow["{i \triangleleft 1}", from=1-4, to=1-6]
	\arrow["\lambda"', from=1-4, to=3-6]
	\arrow["{\mathrm{eval}^c_c}", from=1-6, to=3-6]
	\arrow["{1 \triangleleft \mathrm{eval}_d^{c}}"', from=2-1, to=3-1]
	\arrow["{\mathrm{eval}_e^{d}}"', from=3-1, to=3-3]
\end{tikzcd}\cdpunct[8pt]{.}\]

The converse of this lemma is not true in general. That is, we can't obtain left module structures using left adjoints to internal homs. However, when $\cV$ is a closed symmetric monoidal category, the converse does hold, and the closed module structure lifts from $\nCat{1}$ to $\nCat{\cV}$. In particular, this defines a enriched, tensored, and cotensored category \cite[Prop.~10.1.14]{reihlcht}.

\begin{lemma} \label{lem:poweringrightmod}
	The powering functor $- \pitchfork - \colon \cC \times \cV^{\op} \rightarrow \cC$ defines a canonical right $\cV^{\op}$-module structure. 
\end{lemma}

\begin{proof}
	A left $\cV$-module structure $(\cC,\triangleleft)$ is equivalent to providing a strong monoidal functor $- \triangleleft  \colon \cV \rightarrow \nCat{1}(\cC,\cC)$. The requirement of $x \triangleleft -$ admitting right adjoint is then equivalent to $x \triangleleft -$ admitting a right dual in the endofunctor monoidal category $\nCat{1}(\cC,\cC)$. The right dual construction then defines a \enquote{partial} normal strong monoidal functor $r((\nCat{1}(\cC,\cC))^{\op}) \rightarrow \nCat{1}(\cC,\cC)$ out of the right dualizable endofunctors (see \cite[Sec.~2.10]{etingof2015tenscat}). In particular, we obtain a canonical left $r(\cV^{\op})$-module structure on $\cC$ by postcomposing with this partial right dual functor. Finally, this is equivalent to a canonical right $\cV^{\op}$-module structure on $\cC$ or a canonical right $\cV$-module structure on $\cC^{\op}$. 
\end{proof}

In concrete terms, the right $\cV^{\op}$-module structure on $\cC$ has a right unitor given by the diagram
% https://q.uiver.app/#q=WzAsMyxbMCwwLCJjIFxccGl0Y2hmb3JrIEkiXSxbMiwwLCJJIFxcdHJpYW5nbGVsZWZ0ICggYyBcXHBpdGNoZm9yayBJKSJdLFsyLDIsImMiXSxbMCwxLCJcXGxhbWJkYV57LTF9X3tjIFxccGl0Y2hmb3JrIEl9Il0sWzEsMiwiXFxlcHNpbG9uXntJfV97Y30iXSxbMCwyXV0=
\[\begin{tikzcd}[ampersand replacement=\&,cramped]
	{c \pitchfork I} \&\& {I \triangleleft ( c \pitchfork I)} \\
	\\
	\&\& c
	\arrow["{\lambda^{-1}_{c \pitchfork I}}", from=1-1, to=1-3]
	\arrow[from=1-1, to=3-3]
	\arrow["{\epsilon^{I}_{c}}", from=1-3, to=3-3]
\end{tikzcd}\cdpunct[2pt]{,}\]
where the inverse of this right unitor is given by a morphism $c \rightarrow (c \pitchfork I)$ that is adjoint to the left unitor $I \triangleleft c \rightarrow c$. In particular, this inverse right unitor fits in the diagram
% https://q.uiver.app/#q=WzAsMyxbMCwwLCJJIFxcdHJpYW5nbGVsZWZ0IGMiXSxbMiwwLCJJIFxcdHJpYW5nbGVsZWZ0IChjIFxccGl0Y2hmb3JrIEkpIl0sWzIsMiwiYyJdLFsxLDIsIlxcZXBzaWxvbl5JX3tjfSJdLFswLDIsIlxcbGFtYmRhIiwyXSxbMCwxXV0=
\[\begin{tikzcd}[ampersand replacement=\&,cramped]
	{I \triangleleft c} \&\& {I \triangleleft (c \pitchfork I)} \\
	\\
	\&\& c
	\arrow[from=1-1, to=1-3]
	\arrow["\lambda"', from=1-1, to=3-3]
	\arrow["{\epsilon^I_{c}}", from=1-3, to=3-3]
\end{tikzcd}\cdpunct[2pt]{.}\]

Similarly, the associator of the right module structure on $\cC$ is given by a map $(c \pitchfork x) \pitchfork y  \rightarrow c \pitchfork (x \otimes y) \in \cC$ which fits into the diagram
% https://q.uiver.app/#q=WzAsNSxbMCwwLCIoeCBcXG90aW1lcyB5KSBcXHRyaWFuZ2xlbGVmdCAoKGMgXFxwaXRjaGZvcmsgeCkgXFxwaXRjaGZvcmsgeSkiXSxbMiwwLCIoeCBcXG90aW1lcyB5KSBcXHRyaWFuZ2xlbGVmdCAoYyBcXHBpdGNoZm9yayAoeCBcXG90aW1lcyB5KSkiXSxbMiwzLCJjIl0sWzAsMSwieCBcXHRyaWFuZ2xlbGVmdCAoIHkgXFx0cmlhbmdsZWxlZnQgKChjIFxccGl0Y2hmb3JrIHgpXFxwaXRjaGZvcmsgeSkpIl0sWzAsMywieCBcXHRyaWFuZ2xlbGVmdCAoYyBcXHBpdGNoZm9yayB4KSJdLFswLDFdLFsxLDIsIlxcZXBzaWxvbl57eCBcXG90aW1lcyB5fV97Y30iXSxbMCwzLCJhIl0sWzMsNCwiMSBcXHRyaWFuZ2xlbGVmdCBcXGVwc2lsb25eeV97YyBcXHBpdGNoZm9yayB4fSJdLFs0LDIsIlxcZXBzaWxvbl54X3tjfSIsMl1d
\[\begin{tikzcd}[ampersand replacement=\&,cramped]
	{(x \otimes y) \triangleleft ((c \pitchfork x) \pitchfork y)} \&\& {(x \otimes y) \triangleleft (c \pitchfork (x \otimes y))} \\
	{x \triangleleft ( y \triangleleft ((c \pitchfork x)\pitchfork y))} \\
	\\
	{x \triangleleft (c \pitchfork x)} \&\& c
	\arrow[from=1-1, to=1-3]
	\arrow["a", from=1-1, to=2-1]
	\arrow["{\epsilon^{x \otimes y}_{c}}", from=1-3, to=4-3]
	\arrow["{1 \triangleleft \epsilon^y_{c \pitchfork x}}", from=2-1, to=4-1]
	\arrow["{\epsilon^x_{c}}"', from=4-1, to=4-3]
\end{tikzcd}\cdpunct[8pt]{.}\]

This lemma can be readily dualized to provide a left $\cV$-module structure by tensoring given a right $\cV^{\op}$-module structure by powering. Indeed, a right $\cV^{\op}$-module $\cC$ is equivalent to specifying a left $r(\cV)$-module $\cC^{\op}$. Now, we just apply the previous lemma to this equivalent specification to obtain the desired result. Expectedly, it is also true that an enrichment can be defined solely in terms of the powering without requiring the existence of a tensoring operation.

We can now begin discussing the slicing of closed module structures.

\begin{lemma} \label{lem:sliceovermod}
	Let $\cC$ be a left $\cV$-module, let $(M,\mu,e) \in \cV$ be a monoid object, and let $(Z,\triangleleft_Z) \in \cC$ be a left module object over the monoid $M$. And assume that $\cV$ has pullbacks.
	\begin{enumerate}
		\item If the category $\cV$ is left closed, then the category $\cV/M$ is also left closed using the sliced monoidal product. The same is true for right closure.
		\item If $\cC$ is a left closed left $\cV$-module then $\cC/Z$ is a left closed left-$\cV/M$ module. The same is true for right closure as long as $\cC$ has pullbacks.
	\end{enumerate}
\end{lemma}

\begin{proof} \phantom{}
	\begin{enumerate}
		\item We define the left closed internal hom $\ucV/M(\alpha\colon X \rightarrow M ,\beta\colon Y \rightarrow M)$ and right closed powering $\beta \pitchfork_{/} \alpha$ using the pullback diagrams
		% https://q.uiver.app/#q=WzAsMTAsWzAsMiwiTSJdLFsyLDIsIlxcdWNWKFgsTSkiXSxbMiwwLCJcXHVjVihYLFkpIl0sWzAsMCwiXFx1Y1YvTShcXGFscGhhLFxcYmV0YSkiXSxbMSwyLCJcXHVjVihNLE0pIl0sWzMsMiwiTSJdLFs0LDIsIk0gXFxwaXRjaGZvcmsgTSJdLFs2LDIsIk0gXFxwaXRjaGZvcmsgWCJdLFs2LDAsIlkgXFxwaXRjaGZvcmsgWCJdLFszLDAsIlxcYmV0YSBcXHBpdGNoZm9ya197L30gXFxhbHBoYSJdLFszLDJdLFszLDBdLFswLDQsIlxcbXUiLDJdLFs0LDEsIlxcYWxwaGFeKiIsMl0sWzMsNCwiIiwwLHsic3R5bGUiOnsibmFtZSI6ImNvcm5lci1pbnZlcnNlIn19XSxbOCw3LCJcXGJldGEgXFxwaXRjaGZvcmsgWCJdLFs2LDcsIk0gXFxwaXRjaGZvcmsgXFxhbHBoYSIsMl0sWzUsNiwiXFxtdSIsMl0sWzksOF0sWzksNV0sWzksNywiIiwxLHsic3R5bGUiOnsibmFtZSI6ImNvcm5lci1pbnZlcnNlIn19XSxbMiwxLCJcXGJldGFfKiJdXQ==
		\[\begin{tikzcd}[ampersand replacement=\&,cramped,column sep=scriptsize]
			{\ucV/M(\alpha,\beta)} \&\& {\ucV(X,Y)} \& {\beta \pitchfork_{/} \alpha} \&\&\& {Y \pitchfork X} \\
			\\
			M \& {\ucV(M,M)} \& {\ucV(X,M)} \& M \& {M \pitchfork M} \&\& {M \pitchfork X}
			\arrow[from=1-1, to=1-3]
			\arrow[from=1-1, to=3-1]
			\arrow["\ulcorner"{anchor=center, pos=0.125}, draw=none, from=1-1, to=3-2]
			\arrow["{\beta_*}", from=1-3, to=3-3]
			\arrow[from=1-4, to=1-7]
			\arrow[from=1-4, to=3-4]
			\arrow["\ulcorner"{anchor=center, pos=0.125}, draw=none, from=1-4, to=3-7]
			\arrow["{\beta \pitchfork X}", from=1-7, to=3-7]
			\arrow["\mu"', from=3-1, to=3-2]
			\arrow["{\alpha^*}"', from=3-2, to=3-3]
			\arrow["\mu"', from=3-4, to=3-5]
			\arrow["{M \pitchfork \alpha}"', from=3-5, to=3-7]
		\end{tikzcd}\]
		Thereafter, maps $\gamma \colon Z \rightarrow M$ into $\ucV/M(\alpha,\beta) \rightarrow M$ be can identified with maps $\gamma \otimes_{/M} \alpha \rightarrow \beta$ by using the left closure adjunction of $\cV$. The right closure is same.
		
		\item Let $\alpha: X \rightarrow M$, and $f,g: A,B \rightarrow Z$. The definition of internal hom and powering is similarly given pullback diagrams
		% https://q.uiver.app/#q=WzAsMTAsWzAsMiwiTSJdLFsyLDIsIlxcdWNDKEEsWikiXSxbMiwwLCJcXHVjQyhBLEIpIl0sWzAsMCwiXFx1Y0MvWihmLGcpIl0sWzEsMiwiXFx1Y0MoWixaKSJdLFszLDIsIloiXSxbNCwyLCJaIFxccGl0Y2hmb3JrIE0iXSxbNiwyLCJaIFxccGl0Y2hmb3JrIFgiXSxbNiwwLCJCIFxccGl0Y2hmb3JrIFgiXSxbMywwLCJnIFxccGl0Y2hmb3JrX3svfSBcXGFscGhhIl0sWzMsMl0sWzMsMF0sWzAsNCwiXFx0cmlhbmdsZWxlZnRfe1p9IiwyXSxbNCwxLCJmXioiLDJdLFszLDQsIiIsMCx7InN0eWxlIjp7Im5hbWUiOiJjb3JuZXItaW52ZXJzZSJ9fV0sWzgsNywiZyBcXHBpdGNoZm9yayBYIl0sWzYsNywiWiBcXHBpdGNoZm9yayBcXGFscGhhIiwyXSxbNSw2LCJcXHRyaWFuZ2xlbGVmdF97Wn0iLDJdLFs5LDhdLFs5LDVdLFs5LDcsIiIsMSx7InN0eWxlIjp7Im5hbWUiOiJjb3JuZXItaW52ZXJzZSJ9fV0sWzIsMSwiZ18qIl1d
		\[\begin{tikzcd}[ampersand replacement=\&,cramped]
			{\ucC/Z(f,g)} \&\& {\ucC(A,B)} \& {g \pitchfork_{/} \alpha} \&\&\& {B \pitchfork X} \\
			\\
			M \& {\ucC(Z,Z)} \& {\ucC(A,Z)} \& Z \& {Z \pitchfork M} \&\& {Z \pitchfork X}
			\arrow[from=1-1, to=1-3]
			\arrow[from=1-1, to=3-1]
			\arrow["\ulcorner"{anchor=center, pos=0.125}, draw=none, from=1-1, to=3-2]
			\arrow["{g_*}", from=1-3, to=3-3]
			\arrow[from=1-4, to=1-7]
			\arrow[from=1-4, to=3-4]
			\arrow["\ulcorner"{anchor=center, pos=0.125}, draw=none, from=1-4, to=3-7]
			\arrow["{g \pitchfork X}", from=1-7, to=3-7]
			\arrow["{\triangleleft_{Z}}"', from=3-1, to=3-2]
			\arrow["{f^*}"', from=3-2, to=3-3]
			\arrow["{\triangleleft_{Z}}"', from=3-4, to=3-5]
			\arrow["{Z \pitchfork \alpha}"', from=3-5, to=3-7]
		\end{tikzcd}\]
		The verification of adjunction is similar by using the closure adjunctions of left $\cV$-module $\cC$.
	\end{enumerate}
	
\end{proof}

A dual version of the above lemma is not possible for a comonoid $N$ as there is no transpose for the comultiplication $N \rightarrow N \otimes N$. However, we can still pursue a form of duality, if we instead start with a right $\cV^{\op}$-module $(\cC,\pitchfork)$. This right module structure is meant to encode powering. Thus, we merely ask whether this powering admits tensors using left adjoints or internal homs using mutual right adjoints. More precisely, we use the $\op_1$ duality to encode this as a right $\cV$-module $\cC^{\op}$. Then, equivalently, we ask if this right module is left closed for admitting tensors or right closed for admitting internal homs.

\begin{lemma} \label{lem:sliceundermod}
	Let $\cC^{\op}$ be a right $\cV$-module using $(- \pitchfork -)^{\op}$, let $(M,\mu,e) \in \cV$ be a monoid object, and let $(Z,\triangleright_Z) \in \cC^{\op}$ be a right module object over the monoid $M$ (equivalently, a right comodule object $(Z,\triangleright_Z) \in \cC$ over comonoid $M \in \cV^{\op}$). And assume that $\cV$ has pullbacks.
	\begin{enumerate}
		\item There is a canonical right $\cV/M$-module structure on $\cC^{\op}/Z$, equivalently on $(Z/\cC)^{\op}$. Under the $\op_1$ duality, this is a right $(\cV/M)^{\op}$-module structure on $Z/\cC$ and defines sliced powering.
		\item If $\cC^{\op}$ is a right closed right $\cV$-module then $(Z/\cC)^{\op}$ is a right closed right $\cV/M$-module. This defines a sliced internal hom $Z/\ucC(-,-) \colon (Z/\cC)^{\op} \times (Z/\cC) \rightarrow \cV/M$.
		\item If $\cC^{\op}$ is a left closed right $\cV$ module then $(Z/\cC)^{\op}$ is a left closed right $\cV/M$-module as long as $\cC$ has pushouts (equivalently $\cC^{\op}$ has pullbacks). This defines a sliced left action functor $(- \triangleleft_{/} -) \colon \cV/M \times Z/\cC \rightarrow Z/\cC$.
	\end{enumerate}
\end{lemma}

\begin{proof}
	This whole lemma is merely a direct application of the \cref{lem:sliceovermod} under appropriate dualities. Thus, we only define the functors for clarity. Let $\alpha : X \rightarrow M \in \cV/M$ and $f,g :Z \rightarrow A,B \in Z/\cC$. 
	
	The sliced powering $g^{\alpha}$ is defined by the composite
	% https://q.uiver.app/#q=WzAsNCxbMCwwLCJaIl0sWzIsMSwiWiBcXHBpdGNoZm9yayBNIl0sWzQsMiwiWiBcXHBpdGNoZm9yayBYIl0sWzQsMCwiQiBcXHBpdGNoZm9yayBYIl0sWzAsMywiZyBcXHBpdGNoZm9ya197L30gXFxhbHBoYSJdLFswLDEsIlxcdHJpYW5nbGVyaWdodF97Wn0iLDJdLFsxLDIsIlogXFxwaXRjaGZvcmsgXFxhbHBoYSIsMl0sWzIsMywiZyBcXHBpdGNoZm9yayBYIiwyXV0=
	\[\begin{tikzcd}[ampersand replacement=\&,cramped]
		Z \&\&\&\& {B \pitchfork X} \\
		\&\& {Z \pitchfork M} \\
		\&\&\&\& {Z \pitchfork X}
		\arrow["{g \pitchfork_{/} \alpha}", from=1-1, to=1-5]
		\arrow["{\triangleright_{Z}}"', from=1-1, to=2-3]
		\arrow["{Z \pitchfork \alpha}"', from=2-3, to=3-5]
		\arrow["{g \pitchfork X}"', from=3-5, to=1-5]
	\end{tikzcd}\cdpunct[2pt]{.}\]
	
	
	The internal hom $Z/\ucC(f,g)$ is defined by the pullback
	% https://q.uiver.app/#q=WzAsNSxbMCwwLCJaL1xcdWNDKGYsZykiXSxbMCwyLCJNIl0sWzEsMiwiXFx1Y0MoWixaKSJdLFsyLDIsIlxcdWNDKFosQikiXSxbMiwwLCJcXHVjQyhBLEIpIl0sWzIsMywiZ18qIiwyXSxbNCwzLCJmXioiXSxbMSwyLCJcXHRyaWFuZ2xlcmlnaHRfe1p9IiwyXSxbMCwxXSxbMCw0XSxbMCwzLCIiLDEseyJzdHlsZSI6eyJuYW1lIjoiY29ybmVyLWludmVyc2UifX1dXQ==
	\[\begin{tikzcd}[ampersand replacement=\&,cramped]
		{Z/\ucC(f,g)} \&\& {\ucC(A,B)} \\
		\\
		M \& {\ucC(Z,Z)} \& {\ucC(Z,B)}
		\arrow[from=1-1, to=1-3]
		\arrow[from=1-1, to=3-1]
		\arrow["\ulcorner"{anchor=center, pos=0.125}, draw=none, from=1-1, to=3-3]
		\arrow["{f^*}", from=1-3, to=3-3]
		\arrow["{\triangleright_{Z}}"', from=3-1, to=3-2]
		\arrow["{g_*}"', from=3-2, to=3-3]
	\end{tikzcd}\cdpunct[4pt]{.}\]

	
	The sliced tensoring $\alpha \triangleleft_{Z} f$ is defined by the pushout
	% https://q.uiver.app/#q=WzAsNSxbMiwwLCJYIFxcdHJpYW5nbGVsZWZ0IEEiXSxbMCwwLCJYIFxcdHJpYW5nbGVsZWZ0IFoiXSxbMCwxLCJNIFxcdHJpYW5nbGVsZWZ0IFoiXSxbMCwyLCJaIl0sWzIsMiwiXFxhbHBoYSBcXHRyaWFuZ2xlbGVmdF97L1p9IGYiXSxbMSwyLCJcXGFscGhhIFxcdHJpYW5nbGVsZWZ0IDEiLDJdLFsyLDMsIlxcdHJpYW5nbGVyaWdodF97Wn0iLDJdLFsxLDAsIjEgXFx0cmlhbmdsZWxlZnQgZiJdLFswLDRdLFszLDRdLFs0LDEsIiIsMCx7InN0eWxlIjp7Im5hbWUiOiJjb3JuZXItaW52ZXJzZSJ9fV1d
	\[\begin{tikzcd}[ampersand replacement=\&,cramped]
		{X \triangleleft Z} \&\& {X \triangleleft A} \\
		{M \triangleleft Z} \\
		Z \&\& {\alpha \triangleleft_{/Z} f}
		\arrow["{1 \triangleleft f}", from=1-1, to=1-3]
		\arrow["{\alpha \triangleleft 1}"', from=1-1, to=2-1]
		\arrow[from=1-3, to=3-3]
		\arrow["{\triangleright_{Z}}"', from=2-1, to=3-1]
		\arrow[from=3-1, to=3-3]
		\arrow["\ulcorner"{anchor=center, pos=0.125, rotate=180}, draw=none, from=3-3, to=1-1]
	\end{tikzcd}\cdpunct[2pt]{.}\]
\end{proof}

In the previous two lemma, we obtained the sliced enrichment structures by slicing the tensoring and by slicing the powering, respectively. Therefore, a natural question that remains is if there exists a slicing of the enrichment itself. We motivate this through the following lemma.

\begin{lemma}
	Let $\cC$ be a closed left $\cV$-module and $(M,\mu,e)$ a monoid object in $\cV$. The following data are equivalent.
	\begin{enumerate}
		\item A left module object $(Z, M \triangleleft Z \rightarrow Z)$ over the monoid object $M \in \cV$.
		
		\item A morphism of monoid objects $M \rightarrow \ucC(Z,Z)$.
		
		\item A right comodule object $(Z, Z \rightarrow Z \pitchfork M)$ over the comonoid object $M \in \cV^{\op}$.
	\end{enumerate}
\end{lemma}

\begin{proof}
	We define the monoid morphism $M \rightarrow \ucC(Z,Z)$ as adjoint to the morphism $M \triangleleft Z \rightarrow Z$. In particular, we have a commuting diagram 
	% https://q.uiver.app/#q=WzAsMyxbMCwwLCJNIFxcdHJpYW5nbGVsZWZ0IFoiXSxbMiwyLCJaIl0sWzIsMCwiXFx1Y0MoWixaKSBcXHRyaWFuZ2xlbGVmdCBaIl0sWzAsMV0sWzAsMl0sWzIsMSwiXFxtYXRocm17ZXZhbH1eWl9aIl1d
	\[\begin{tikzcd}[ampersand replacement=\&,cramped]
		{M \triangleleft Z} \&\& {\ucC(Z,Z) \triangleleft Z} \\
		\\
		\&\& Z
		\arrow[from=1-1, to=1-3]
		\arrow[from=1-1, to=3-3]
		\arrow["{\mathrm{eval}^Z_Z}", from=1-3, to=3-3]
	\end{tikzcd}\cdpunct[2pt]{.}\]
	Similarly, we define the comodule map as an adjoint to the morphism $M \triangleleft Z \rightarrow Z$, which fits into the following commuting diagram
	% https://q.uiver.app/#q=WzAsMyxbMCwwLCJNIFxcdHJpYW5nbGVsZWZ0IFoiXSxbMiwyLCJaIl0sWzIsMCwiTSBcXHRyaWFuZ2xlbGVmdCAoWiBcXHBpdGNoZm9yayBNKSJdLFsyLDEsIlxcZXBzaWxvbl5NX3tafSJdLFswLDJdLFswLDFdXQ==
	\[\begin{tikzcd}[ampersand replacement=\&,cramped]
		{M \triangleleft Z} \&\& {M \triangleleft (Z \pitchfork M)} \\
		\\
		\&\& Z
		\arrow[from=1-1, to=1-3]
		\arrow[from=1-1, to=3-3]
		\arrow["{\epsilon^M_{Z}}", from=1-3, to=3-3]
	\end{tikzcd}\cdpunct[2pt]{.}\]
	
	The rest is simply an exercise in taking transposes of the defining axioms of the left module structure of $(Z, M \triangleleft Z \rightarrow Z)$ and using the above two diagrams to simplify.
\end{proof}

In particular, this suggests that we should be able to perform slice construction on arbitrary enriched categories $\cC$ over some base $\cV$, as long as we have a morphism of monoids $M \rightarrow \ucC(Z,Z)$. This is indeed true.

\begin{proposition}
	Consider a monoidal category $\cV$ with pullbacks. Let $\ucC$ be any enriched category and assume that there exists a morphism of monoid objects $r \colon M \rightarrow \ucC(Z,Z)$ for some $Z \in \ucC$. There exists canonical slice enriched categories $Z/\ucC$ and $\ucC/Z$ enriched over the slice monoidal category $\cV/M$. Furthermore, the underlying category of both $\ucC/Z$ and $Z/\ucC$ over $\cV/M$ (and not $\cV$) is isomorphic to the corresponding slices categories $\cC/Z$ and $Z/\cC$. 
\end{proposition}

\begin{proof}
	Let the maps $f_*$ and $f^*$ be formed by taking the unenriched representables. We only define the compositor and the unitor for the case of $\ucC/Z$ enrichment over $\cV/M$. The case for $Z/\ucC$ is identical. 
	
	The sliced internal homs are defined by the pullbacks % https://q.uiver.app/#q=WzAsNSxbMCwyLCJNIl0sWzIsMiwiXFx1Y0MoQSxaKSJdLFsyLDAsIlxcdWNDKEEsQikiXSxbMCwwLCJcXHVjQy9aKGYsZykiXSxbMSwyLCJcXHVjQyhaLFopIl0sWzMsMl0sWzMsMF0sWzAsNCwiciIsMl0sWzQsMSwiZl4qIiwyXSxbMyw0LCIiLDAseyJzdHlsZSI6eyJuYW1lIjoiY29ybmVyLWludmVyc2UifX1dLFsyLDEsImdfKiJdXQ==
	\[\begin{tikzcd}[ampersand replacement=\&,cramped]
		{\ucC/Z(f,g)} \&\& {\ucC(A,B)} \\
		\\
		M \& {\ucC(Z,Z)} \& {\ucC(A,Z)}
		\arrow[from=1-1, to=1-3]
		\arrow[from=1-1, to=3-1]
		\arrow["\ulcorner"{anchor=center, pos=0.125}, draw=none, from=1-1, to=3-2]
		\arrow["{g_*}", from=1-3, to=3-3]
		\arrow["r"', from=3-1, to=3-2]
		\arrow["{f^*}"', from=3-2, to=3-3]
	\end{tikzcd}\cdpunct[8pt]{.}\]
	
	The unitors are defined by the unique universal maps $i \colon I \rightarrow \ucC(A,A)$ as given in the diagrams
	% https://q.uiver.app/#q=WzAsNixbMSwzLCJNIl0sWzMsMywiXFx1Y0MoQSxaKSJdLFszLDEsIlxcdWNDKEEsQSkiXSxbMSwxLCJcXHVjQy9aKGYsZikiXSxbMiwzLCJcXHVjQyhaLFopIl0sWzAsMCwiSSJdLFszLDJdLFszLDBdLFswLDQsInIiLDJdLFs0LDEsImZeKiIsMl0sWzMsNCwiIiwwLHsic3R5bGUiOnsibmFtZSI6ImNvcm5lci1pbnZlcnNlIn19XSxbMiwxLCJmXyoiXSxbNSwwLCJlIiwyLHsiY3VydmUiOjJ9XSxbNSwyLCIxX3tBfSIsMCx7ImN1cnZlIjotMn1dLFs1LDMsImkiLDFdXQ==
	\[\begin{tikzcd}[ampersand replacement=\&,cramped]
		I \&\&\& \\
		\& {\ucC/Z(f,f)} \&\& {\ucC(A,A)} \\
		\\
		\& M \& {\ucC(Z,Z)} \& {\ucC(A,Z)}
		\arrow["i"{description}, from=1-1, to=2-2]
		\arrow["{1_{A}}", curve={height=-12pt}, from=1-1, to=2-4]
		\arrow["e"', curve={height=12pt}, from=1-1, to=4-2]
		\arrow[from=2-2, to=2-4]
		\arrow[from=2-2, to=4-2]
		\arrow["\ulcorner"{anchor=center, pos=0.125}, draw=none, from=2-2, to=4-3]
		\arrow["{f_*}", from=2-4, to=4-4]
		\arrow["r"', from=4-2, to=4-3]
		\arrow["{f^*}"', from=4-3, to=4-4]
	\end{tikzcd}\cdpunct[8pt]{,}\]
	where the outer rectangle commutes because $r$ is a morphism of monoids.
	
	Finally, the compositors $\ucC/Z(g,h) \otimes \ucC/Z(f,g) \rightarrow \ucC/Z(f,h)$ are the unique universal maps induced by the following commuting diagrams
	% https://q.uiver.app/#q=WzAsMTYsWzAsMCwiXFx1Y0MvWihnLGgpIFxcb3RpbWVzIFxcdWNDL1ooZixnKSJdLFsxLDAsIlxcdWNDL1ooZyxoKSBcXG90aW1lcyBcXHVjQyhBLEIpIl0sWzAsMiwiTSBcXG90aW1lcyBcXHVjQy9aKGYsZykiXSxbMiwwLCJcXHVjQyhCLEMpIFxcb3RpbWVzIFxcdWNDKEEsQikiXSxbMywwLCJcXHVjQyhBLEMpIl0sWzAsNCwiTSBcXG90aW1lcyBNIl0sWzAsNiwiTSJdLFszLDYsIlxcdWNDKFosWikiXSxbMywyLCJcXHVjQyhBLFopIl0sWzIsMSwiXFx1Y0MoQixaKSBcXG90aW1lcyBcXHVjQyhBLEIpIl0sWzIsMiwiXFx1Y0MoWixaKSBcXG90aW1lcyBcXHVjQyhBLEIpIl0sWzEsMiwiTSBcXG90aW1lcyBcXHVjQyhBLEIpIl0sWzEsNCwiTSBcXG90aW1lcyBcXHVjQyhaLFopIl0sWzIsNCwiXFx1Y0MoWixaKSBcXG90aW1lcyBcXHVjQyhaLFopIl0sWzIsMywiXFx1Y0MoWixaKSBcXG90aW1lcyBcXHVjQyhBLFopIl0sWzEsMywiTSBcXG90aW1lcyBcXHVjQyhBLFopIl0sWzAsMl0sWzAsMV0sWzEsM10sWzMsNCwiXFxjaXJjIl0sWzIsNV0sWzUsNiwiXFxtdSIsMl0sWzYsNywiciIsMl0sWzcsOCwiZl4qIiwyXSxbNCw4LCJoXyoiXSxbMyw5LCJoXyogXFxvdGltZXMgMSIsMl0sWzksOCwiXFxjaXJjIl0sWzEwLDksImdeKiBcXG90aW1lcyAxIl0sWzIsMTFdLFsxMSwxMCwiciBcXG90aW1lcyAxIl0sWzEyLDEzLCJyIFxcb3RpbWVzIDEiLDJdLFsxMywxNCwiMSBcXG90aW1lcyBmXioiXSxbNSwxMiwiMSBcXG90aW1lcyByIiwyXSxbMTMsNywiXFxjaXJjIiwyXSxbMTIsMTUsIjEgXFxvdGltZXMgZl4qIl0sWzE1LDE0LCJyIFxcb3RpbWVzIDEiLDJdLFsxMSwxNSwiMSBcXG90aW1lcyBnXyoiLDJdLFsxMCwxNCwiMSBcXG90aW1lcyBnXyoiLDJdLFsxNCw4LCJcXGNpcmMiLDJdLFsxLDExXV0=
	\[\begin{tikzcd}[ampersand replacement=\&,cramped,column sep=small]
		{\ucC/Z(g,h) \otimes \ucC/Z(f,g)} \& {\ucC/Z(g,h) \otimes \ucC(A,B)} \& {\ucC(B,C) \otimes \ucC(A,B)} \& {\ucC(A,C)} \\
		\&\& {\ucC(B,Z) \otimes \ucC(A,B)} \\
		{M \otimes \ucC/Z(f,g)} \& {M \otimes \ucC(A,B)} \& {\ucC(Z,Z) \otimes \ucC(A,B)} \& {\ucC(A,Z)} \\
		\& {M \otimes \ucC(A,Z)} \& {\ucC(Z,Z) \otimes \ucC(A,Z)} \\
		{M \otimes M} \& {M \otimes \ucC(Z,Z)} \& {\ucC(Z,Z) \otimes \ucC(Z,Z)} \\
		\\
		M \&\&\& {\ucC(Z,Z)}
		\arrow[from=1-1, to=1-2]
		\arrow[from=1-1, to=3-1]
		\arrow[from=1-2, to=1-3]
		\arrow[from=1-2, to=3-2]
		\arrow["\circ", from=1-3, to=1-4]
		\arrow["{h_* \otimes 1}"', from=1-3, to=2-3]
		\arrow["{h_*}", from=1-4, to=3-4]
		\arrow["\circ", from=2-3, to=3-4]
		\arrow[from=3-1, to=3-2]
		\arrow[from=3-1, to=5-1]
		\arrow["{r \otimes 1}", from=3-2, to=3-3]
		\arrow["{1 \otimes g_*}"', from=3-2, to=4-2]
		\arrow["{g^* \otimes 1}", from=3-3, to=2-3]
		\arrow["{1 \otimes g_*}"', from=3-3, to=4-3]
		\arrow["{r \otimes 1}"', from=4-2, to=4-3]
		\arrow["\circ"', from=4-3, to=3-4]
		\arrow["{1 \otimes r}"', from=5-1, to=5-2]
		\arrow["\mu"', from=5-1, to=7-1]
		\arrow["{1 \otimes f^*}", from=5-2, to=4-2]
		\arrow["{r \otimes 1}"', from=5-2, to=5-3]
		\arrow["{1 \otimes f^*}", from=5-3, to=4-3]
		\arrow["\circ"', from=5-3, to=7-4]
		\arrow["r"', from=7-1, to=7-4]
		\arrow["{f^*}"', from=7-4, to=3-4]
	\end{tikzcd}\cdpunct[8pt]{.}\]
	
	For the associativity and unitality, first note that $U \colon \cV/M \rightarrow \cV$ is a faithful strict monoidal functor. Thus, $f = g$ in $\cV/M$ iff $U(f)=U(g)$ in $\cV$. Second, universal maps defined through pullbacks are uniquely characterized by the two projection maps. 
	
	Using these two facts we shall, for instance, demonstrate the left unitality. The associativity is quite similar but only technically more involved. Now, note the commuting diagram
	% https://q.uiver.app/#q=WzAsOSxbMCwwLCJJIFxcb3RpbWVzIFxcdWNDL1ooZixnKSJdLFswLDUsIlxcdWNDL1ooZyxnKSBcXG90aW1lcyBcXHVjQy9aKGYsZykiXSxbMiw1LCJcXHVjQy9aKGYsZykiXSxbMyw1LCJcXHVjQyhBLEIpIl0sWzEsMiwiXFx1Y0MoQixCKSBcXG90aW1lcyBcXHVjQy9aKGYsZykiXSxbMiw0LCJcXHVjQyhCLEIpIFxcb3RpbWVzIFxcdWNDKEEsQikiXSxbMiwxLCJJIFxcb3RpbWVzIFxcdWNDKEEsQikiXSxbMywwLCJcXHVjQy9aKGYsZykiXSxbMSw0LCJcXHVjQy9aKGcsZykgXFxvdGltZXMgXFx1Y0MoQSxCKSJdLFsxLDJdLFswLDFdLFsyLDNdLFsxLDRdLFs0LDVdLFs1LDMsIlxcY2lyYyIsMl0sWzAsNCwiMV9CIFxcb3RpbWVzIDEiXSxbMCw2XSxbNiw1LCIxX0IgXFxvdGltZXMgMSIsMl0sWzYsMywiXFxsYW1iZGEiXSxbMCw3LCJcXGxhbWJkYSJdLFs3LDNdLFsxLDhdLFs4LDVdXQ==
	\[\begin{tikzcd}[ampersand replacement=\&,cramped]
		{I \otimes \ucC/Z(f,g)} \&\&\& {\ucC/Z(f,g)} \\
		\&\& {I \otimes \ucC(A,B)} \\
		\& {\ucC(B,B) \otimes \ucC/Z(f,g)} \\
		\\
		\& {\ucC/Z(g,g) \otimes \ucC(A,B)} \& {\ucC(B,B) \otimes \ucC(A,B)} \\
		{\ucC/Z(g,g) \otimes \ucC/Z(f,g)} \&\& {\ucC/Z(f,g)} \& {\ucC(A,B)}
		\arrow["\lambda", from=1-1, to=1-4]
		\arrow[from=1-1, to=2-3]
		\arrow["{1_B \otimes 1}", from=1-1, to=3-2]
		\arrow[from=1-1, to=6-1]
		\arrow[from=1-4, to=6-4]
		\arrow["{1_B \otimes 1}"', from=2-3, to=5-3]
		\arrow["\lambda", from=2-3, to=6-4]
		\arrow[from=3-2, to=5-3]
		\arrow[from=5-2, to=5-3]
		\arrow["\circ"', from=5-3, to=6-4]
		\arrow[from=6-1, to=3-2]
		\arrow[from=6-1, to=5-2]
		\arrow[from=6-1, to=6-3]
		\arrow[from=6-3, to=6-4]
	\end{tikzcd}\]
	implies that the left unitality holds under taking postcomposite with one of the projection maps. Similarly, we have the diagram
	% https://q.uiver.app/#q=WzAsOCxbNCwwLCJcXHVjQy9aKGcsZykgIFxcb3RpbWVzIFxcdWNDL1ooZixnKSJdLFswLDAsIkkgXFxvdGltZXMgXFx1Y0MvWihmLGcpIl0sWzQsMiwiXFx1Y0MvWihmLGcpIl0sWzIsNSwiTSJdLFszLDEsIk0gXFxvdGltZXMgXFx1Y0MvWihmLGcpIl0sWzMsMiwiTSBcXG90aW1lcyBNIl0sWzIsMiwiSSBcXG90aW1lcyBNIl0sWzAsMiwiXFx1Y0MvWihmLGcpIl0sWzEsMF0sWzAsMl0sWzIsM10sWzAsNF0sWzQsNV0sWzUsM10sWzEsNCwiZSBcXG90aW1lcyAxIiwyXSxbMSw2XSxbNiw1LCJlIFxcb3RpbWVzIDEiLDJdLFs2LDMsIlxcbGFtYmRhIiwyXSxbMSw3LCJcXGxhbWJkYSIsMl0sWzcsM11d
	\[\begin{tikzcd}[ampersand replacement=\&,cramped]
		{I \otimes \ucC/Z(f,g)} \&\&\&\& {\ucC/Z(g,g)  \otimes \ucC/Z(f,g)} \\
		\&\&\& {M \otimes \ucC/Z(f,g)} \\
		{\ucC/Z(f,g)} \&\& {I \otimes M} \& {M \otimes M} \& {\ucC/Z(f,g)} \\
		\\
		\\
		\&\& M
		\arrow[from=1-1, to=1-5]
		\arrow["{e \otimes 1}"', from=1-1, to=2-4]
		\arrow["\lambda"', from=1-1, to=3-1]
		\arrow[from=1-1, to=3-3]
		\arrow[from=1-5, to=2-4]
		\arrow[from=1-5, to=3-5]
		\arrow[from=2-4, to=3-4]
		\arrow[from=3-1, to=6-3]
		\arrow["{e \otimes 1}"', from=3-3, to=3-4]
		\arrow["\lambda"', from=3-3, to=6-3]
		\arrow[from=3-4, to=6-3]
		\arrow[from=3-5, to=6-3]
	\end{tikzcd}\]
	for compatibility with the other projection map. Finally, these two commuting diagrams jointly imply the desired commuting diagram 
	% https://q.uiver.app/#q=WzAsMyxbMSwwLCJcXHVjQy9aKGcsZykgIFxcb3RpbWVzIFxcdWNDL1ooZixnKSJdLFswLDIsIkkgXFxvdGltZXMgXFx1Y0MvWihmLGcpIl0sWzIsMiwiXFx1Y0MvWihmLGcpIl0sWzEsMF0sWzAsMl0sWzEsMiwiXFxsYW1iZGEiLDJdXQ==
	\[\begin{tikzcd}[ampersand replacement=\&,cramped]
		\& {\ucC/Z(g,g)  \otimes \ucC/Z(f,g)} \& \\
		\\
		{I \otimes \ucC/Z(f,g)} \&\& {\ucC/Z(f,g)}
		\arrow[from=1-2, to=3-3]
		\arrow[from=3-1, to=1-2]
		\arrow["\lambda"', from=3-1, to=3-3]
	\end{tikzcd}\]
	for left unitality. As already mentioned, the verification for associativity and the underlying category is mostly similar.
	
	
\end{proof}

Let $\ucC$ be any $\cV$ enriched category. There is a canonical monoid structure on the monoidal unit $I$ of $\cV$ given by $\mu : I \otimes I \xrightarrow{\lambda_I = \rho_I} I$ and $\eta : I \xrightarrow{1_I} I$. In particular, $\cV/I$ is a semicartesian monoidal category. That is, the induced monoidal unit $I \rightarrow I$ is a terminal object of $\cV/I$. Furthermore, for any $Z \in \cC$, the unitor map $I \rightarrow \ucC(Z,Z)$ is a morphism of monoid objects. This readily provides the following corollary.

\begin{corollary}
	There exists canonical enriched slice categories $\ucC/Z$, $Z/\ucC$ over $\cV/I$ for every $Z \in \cC$.
\end{corollary}

\begin{corollary}
	Let $\cV$ be a closed symmetric monooidal category, and let $\cC$ be enriched, tensored, and cotensored over $\cV$. Assume the relevant pullbacks and pushouts exist. Then, both the categories $\ucC/Z$ and $Z/\ucC$ are enriched, tensored, and cotensored over $\ucV/I$ for every $Z \in \ucC$. More generally, this is also true for a commutative monoid object $M$. 
\end{corollary}

\subsection{Model structures and slices} \label{sub:slicecatmod}

Monoidal model categories and enriched model categories are an important tool in abstract homotopy theory. Their defining axioms are closely linked with the two variable adjunction defining the closure of left $\cV$-modules. Hence, we can also analyze the slicing of enriched model structures.

\begin{comment}
	% https://q.uiver.app/#q=WzAsNSxbMCwwLCJYXzEgXFxvdGltZXMgWV8xIl0sWzIsMywiWF8yIFxcb3RpbWVzIFlfMiJdLFsyLDAsIlhfMSBcXG90aW1lcyBZXzIiXSxbMCwzLCJYXzIgXFxvdGltZXMgWV8xIl0sWzEsMiwiWF8xIFxcb3RpbWVzIFlfMiBcXGRpc3BsYXlzdHlsZVxcY29wcm9kX3tYXzEgXFxvdGltZXMgWV8xfSBYXzIgXFxvdGltZXMgWV8yIl0sWzMsMSwiMSBcXG90aW1lcyBcXGJldGEiLDJdLFsyLDEsIlxcYWxwaGEgXFxvdGltZXMgMSJdLFswLDMsIlxcYWxwaGEgXFxvdGltZXMgMSIsMl0sWzAsMiwiMSBcXG90aW1lcyBcXGJldGEiXSxbNCwxLCJcXGFscGhhIFxcYmFye1xcb3RpbWVzfSBcXGJldGEiXSxbMyw0XSxbMiw0XSxbNCwwLCIiLDEseyJzdHlsZSI6eyJuYW1lIjoiY29ybmVyLWludmVyc2UifX1dXQ==
	\[\begin{tikzcd}
		{X_1 \otimes Y_1} && {X_1 \otimes Y_2} \\
		\\
		& {X_1 \otimes Y_2 \displaystyle\coprod_{X_1 \otimes Y_1} X_2 \otimes Y_2} \\
		{X_2 \otimes Y_1} && {X_2 \otimes Y_2}
		\arrow["{1 \otimes \beta}", from=1-1, to=1-3]
		\arrow["{\alpha \otimes 1}"', from=1-1, to=4-1]
		\arrow[from=1-3, to=3-2]
		\arrow["{\alpha \otimes 1}", from=1-3, to=4-3]
		\arrow["\ulcorner"{anchor=center, pos=0.125, rotate=180}, draw=none, from=3-2, to=1-1]
		\arrow["{\alpha \bar{\otimes} \beta}", from=3-2, to=4-3]
		\arrow[from=4-1, to=3-2]
		\arrow["{1 \otimes \beta}"', from=4-1, to=4-3]
	\end{tikzcd}\]
\end{comment}


\begin{definition}
	A monoidal model category is a \emph{closed} monoidal category  $(\cV,\otimes, I)$ together with a model structure on its underlying category $\cV$ that satisfies:
	\begin{enumerate}
		\item Unit axiom : Let $QI$ be any cofibrant replacement of $I$. The induced maps $QI \otimes X \rightarrow I \otimes X$ and $X \otimes QI \rightarrow X \otimes I$ are weak equivalences, whenever $X$ is cofibrant in $\cV$. 
		\item Pushout product axiom : Let $\alpha : X_1 \rightarrow Y_1$ and $\beta: X_2 \rightarrow Y_2$ be two arbitrary maps in  $\cV$. The pushout-product map $\alpha \bar{\otimes} \beta$ is a cofibration whenever both $\alpha$ and $\beta$ are cofibrations. Furthermore, if one of $\alpha$ or $\beta$ is also a weak equivalence then so is $\alpha \bar{\otimes} \beta$.
	\end{enumerate}
\end{definition} 

\begin{remark}
	The closedness assumption of the monoidal structure on $\cV$ is not entirely redundant. One of its important consequences is that the monoidal product $\otimes$ preserves the initial object in both of its arguments. In fact, it is sufficient to assume only this additional property for constructing the derived monoidal product using a left deformation. See \cite[Sec.~10.2]{reihlcht} or \cite[Prop.~4.3.1]{hovey2007modelcat} for more details.
\end{remark}

We now verify that this extends to a natural monoidal model structure for $\cV/M$. 


\begin{proposition}
	Let $(\cV,\otimes, I)$ be a monoidal model category, $(M, \mu, \eta)$ be a monoid object inside of this monoidal category. Then, $(\cV/M, \otimes_{/M}, I_{/M})$ is also a monoidal model category equipped with the slice model structure. 
\end{proposition}

\begin{proof}
	The slice model structure is defined by the property that all of its cofibrations, fibrations, and weak equivalences are created under the forgetful functor $U : \cV/M \rightarrow \cV$. 
	\begin{enumerate}
		\item Unit Axiom : Any cofibrant replacement of $\eta : I \rightarrow M$ must also be created under the forgetful functor $U$. In particular, it must be of the form $QI \rightarrow I \rightarrow M$, where we take the functorial cofibrant replacement in $\cV$ itself. Now let $X \rightarrow M \in \cV/M$ be any cofibrant object. We obtain that the following diagram which demonstrates the required unit axiom.
		% https://q.uiver.app/#q=WzAsNSxbMCwwLCJRSSBcXG90aW1lcyBYIl0sWzIsMCwiSSBcXG90aW1lcyBYIl0sWzEsMiwiTSBcXG90aW1lcyBYIl0sWzEsMywiTSBcXG90aW1lcyBNIl0sWzEsNCwiTSJdLFswLDFdLFswLDJdLFsxLDJdLFsyLDNdLFszLDRdXQ==
		\[\begin{tikzcd}
			{QI \otimes X} && {I \otimes X} \\
			\\
			& {M \otimes X} \\
			& {M \otimes M} \\
			& M
			\arrow[from=1-1, to=1-3]
			\arrow[from=1-1, to=3-2]
			\arrow[from=1-3, to=3-2]
			\arrow[from=3-2, to=4-2]
			\arrow[from=4-2, to=5-2]
		\end{tikzcd}\]
		\item Pushout product axiom : The functor $U$ is strict monoidal and creates all colimits. Thus, the pushout product axiom of $\cV$ lifts to that of $\cV/M$ directly.  
	\end{enumerate}
\end{proof}

Having obtained the monoidal model category structure on $\cV/M$, we can now analyze the slicing of closed left $\cV$-modules under consideration of model structures. In particular, we start with a more general definition of a $\cV$-module model category instead of a enriched model category. 

\begin{comment}
	% https://q.uiver.app/#q=WzAsNSxbMCwwLCJYIFxcdHJpYW5nbGVsZWZ0IEEiXSxbMiwzLCJZIFxcdHJpYW5nbGVsZWZ0IEIiXSxbMiwwLCJYIFxcdHJpYW5nbGVsZWZ0IEIiXSxbMCwzLCJZIFxcdHJpYW5nbGVsZWZ0IEEiXSxbMSwyLCJYIFxcdHJpYW5nbGVsZWZ0IEIgXFxkaXNwbGF5c3R5bGVcXGNvcHJvZF97WCBcXHRyaWFuZ2xlbGVmdCBBfSBZIFxcdHJpYW5nbGVsZWZ0IEEiXSxbMywxLCIxIFxcdHJpYW5nbGVsZWZ0IGYiLDJdLFsyLDEsIlxcYWxwaGEgXFx0cmlhbmdsZWxlZnQgMSJdLFswLDMsIlxcYWxwaGEgXFx0cmlhbmdsZWxlZnQgMSIsMl0sWzAsMiwiMSBcXHRyaWFuZ2xlbGVmdCBmIl0sWzQsMSwiXFxhbHBoYSBcXGJhcntcXHRyaWFuZ2xlbGVmdH0gZiJdLFszLDRdLFsyLDRdLFs0LDAsIiIsMSx7InN0eWxlIjp7Im5hbWUiOiJjb3JuZXItaW52ZXJzZSJ9fV1d
	\[\begin{tikzcd}[ampersand replacement=\&,cramped]
		{X \triangleleft A} \&\& {X \triangleleft B} \\
		\\
		\& {X \triangleleft B \displaystyle\coprod_{X \triangleleft A} Y \triangleleft A} \\
		{Y \triangleleft A} \&\& {Y \triangleleft B}
		\arrow["{1 \triangleleft f}", from=1-1, to=1-3]
		\arrow["{\alpha \triangleleft 1}"', from=1-1, to=4-1]
		\arrow[from=1-3, to=3-2]
		\arrow["{\alpha \triangleleft 1}", from=1-3, to=4-3]
		\arrow["\ulcorner"{anchor=center, pos=0.125, rotate=180}, draw=none, from=3-2, to=1-1]
		\arrow["{\alpha \bar{\triangleleft} f}", from=3-2, to=4-3]
		\arrow[from=4-1, to=3-2]
		\arrow["{1 \triangleleft f}"', from=4-1, to=4-3]
	\end{tikzcd}\]
\end{comment}

\begin{definition}
	Let $(\cV,\otimes, I)$ be a monoidal model category. A left $\cV$-module model  category $\cC$ is a \emph{closed} left $\cV$-module $(\cC, \triangleleft)$ together with a model structure on $\cC$ that satisfies:
	\begin{enumerate}
		\item Tensoring unit axiom : Let $QI$ be any cofibrant replacement of $I$. The induced map $QI \triangleleft A \rightarrow I \triangleleft A$ is a weak equivalences, whenever $A$ is cofibrant in $\cC$. 
		\item Pushout product axiom : Let $\alpha : X \rightarrow Y \in \cV$ and $f : A \rightarrow B \in \cC$ be two arbitrary maps. The pushout-product map $\alpha \bar{\triangleleft} f$ is a cofibration, whenever both $\alpha$ and $f$ are cofibrations. Furthermore, if one of $\alpha$ or $f$ is a also weak equivalence  then so is $\alpha \bar{\triangleleft} f$.
	\end{enumerate}
\end{definition}

Motivated by the \cref{lem:sliceundermod}, we also have a dual definition. 

\begin{definition}
	Let $(\cV,\otimes, I)$ be a monoidal model category. A right $\cV^{\op}$-module model category $\cC$ is a \emph{closed} right $\cV$-module $(\cC^{\op}, \pitchfork^{\op})$ together with a model structure on $\cC$ that satisfies:
	\begin{enumerate}
		\item Powering unit axiom : Let $QI$ be any cofibrant replacement of $I$. The induced map $A \pitchfork I \rightarrow A \pitchfork QI$ is a weak equivalences, whenever $A$ is fibrant in $\cC$. 
		\item Pullback power axiom : Let $\alpha : X \rightarrow Y \in \cV$ and $f : A \rightarrow B \in \cC$ be two arbitrary maps. The pullback-power map $f \bar{\pitchfork} \alpha$ is a fibration, whenever $\alpha$ is a cofibration and $f$ is a fibrations. Furthermore, if one of $\alpha$ or $f$ is a also weak equivalence then so is $f \bar{\pitchfork} \alpha$.
	\end{enumerate}
\end{definition}

\begin{remark}
	Under the application of both $\op_1$ and $r$ dualities, the above definition specify precisely the tensoring unit axiom and the pushout-product axioms for the corresponding left $r(\cV)$-module structure on $\cC^{\op}$. In particular, the above definition is equivalent to specifying a left $r(\cV)$-module model category $\cC^{\op}$. The primary reason for having distinct axioms is that while $r(\cV)$ is again a monoidal model category, $\op_1(\cV)$ is not, in general.
\end{remark}

As before, there is also a correspondence between left $\cV$-module model categories and right $\cV^{\op}$-module model categories. This primarily follows from \cref{lem:poweringrightmod} together with its dual and \cite[Lemmas~4.2.2 and 4.2.7]{hovey2007modelcat}. 

\begin{theorem}
	Let $(\cV,\otimes, I)$ be a monoidal model category and let $(\cC,\triangleleft)$ be a left $\cV$-module model category. Let $(M, \mu, \eta)$ be a monoid object in $\cV$ and let $(Z,\lambda)$ be a left $M$-module in $\cC$. Then, it follows that $(\cC/Z,\triangleleft_{/Z})$ is a $\cV/M$-model category. 
\end{theorem}

\begin{proof}
	The slice model structure on $\cC/Z$ is again defined by the property that all of the cofibrations, fibrations, and weak equivalences are created under the forgetful functor $U : \cC/Z \rightarrow \cC$. 
	
	The rest of the proof is also similar to the case of a monoidal model category. Indeed, the idea of forgetful functors being \enquote{strict monoidal} is now encoded by the commuting diagram
	% https://q.uiver.app/#q=WzAsNCxbMCwwLCJcXGNWL00gXFx0aW1lcyBcXGNDL1oiXSxbMiwwLCJcXGNDL1oiXSxbMCwyLCJcXGNWIFxcdGltZXMgXFxjQyJdLFsyLDIsIlxcY0MiXSxbMCwyLCJVIFxcdGltZXMgVSIsMl0sWzEsMywiVSJdLFswLDEsIlxcdHJpYW5nbGVsZWZ0X3svWn0iXSxbMiwzLCJcXHRyaWFuZ2xlbGVmdCIsMl1d
	\[\begin{tikzcd}[ampersand replacement=\&,cramped]
		{\cV/M \times \cC/Z} \&\& {\cC/Z} \\
		\\
		{\cV \times \cC} \&\& \cC
		\arrow["{\triangleleft_{/Z}}", from=1-1, to=1-3]
		\arrow["{U \times U}"', from=1-1, to=3-1]
		\arrow["U", from=1-3, to=3-3]
		\arrow["\triangleleft"', from=3-1, to=3-3]
	\end{tikzcd}\cdpunct[8pt]{.}\]
\end{proof}

It is not possible to use the previous technique in the case of under categories $Z/\cC$. In particular, we only have the commuting diagram
% https://q.uiver.app/#q=WzAsNCxbMCwwLCJcXGNWL00gXFx0aW1lcyBaL1xcY0MiXSxbMiwwLCJaL1xcY0MiXSxbMCwyLCJcXGNWIFxcdGltZXMgXFxjQyJdLFsyLDIsIlxcY0MiXSxbMCwyLCJVIFxcdGltZXMgVSIsMl0sWzEsMywiVSJdLFswLDEsIlxcdHJpYW5nbGVsZWZ0X3svWn0iXSxbMiwzLCJcXHRyaWFuZ2xlbGVmdCIsMl0sWzIsMSwibCIsMix7InNob3J0ZW4iOnsic291cmNlIjoyMCwidGFyZ2V0IjoyMH0sImxldmVsIjoyfV1d
\[\begin{tikzcd}[ampersand replacement=\&,cramped]
	{\cV/M \times Z/\cC} \&\& {Z/\cC} \\
	\\
	{\cV \times \cC} \&\& \cC
	\arrow["{\triangleleft_{/Z}}", from=1-1, to=1-3]
	\arrow["{U \times U}"', from=1-1, to=3-1]
	\arrow["U", from=1-3, to=3-3]
	\arrow["l"', between={0.2}{0.8}, Rightarrow, nfold, from=3-1, to=1-3]
	\arrow["\triangleleft"', from=3-1, to=3-3]
\end{tikzcd}\]
up to a $2$-cell $l\colon X \triangleleft A \rightarrow \alpha \triangleleft_{/Z} f$, encoding a \enquote{lax monoidal} functor.

We can attempt to remedy this situation by demanding that the canonical map $l$ be a natural weak equivalence. While this may potentially solve the issue, it is not a desirable assumption in practice. Instead, as before, the solution is to dualize the results with right $\cV^{\op}$-module structures.

\begin{theorem}
	Let $(\cV,\otimes, I)$ be a monoidal model category and let $(\cC,\pitchfork)$ be a right $\cV^{\op}$-module model category. Let $(M, \mu, \eta)$ be a monoid object in $\cV$ and let $(Z,\rho)$ be a right $M$-module in $\cC^{\op}$. Then, it follows that $(Z/\cC, \pitchfork_{/Z})$ is a right $(\cV/M)^{\op}$-module model category. Consequentially, this automatically implies $(Z/\cC, \triangleleft_{/Z})$ is a left $\cV/M$-module model category.
\end{theorem}

\begin{proof}
	The idea of forgetful functors being \enquote{strict monoidal} is now encoded by the following commuting diagram 
	% https://q.uiver.app/#q=WzAsNCxbMCwwLCJaL1xcY0MgXFx0aW1lcyAoXFxjVi9NKV57XFxvcH0iXSxbMiwwLCJaL1xcY0MiXSxbMiwyLCJcXGNDIl0sWzAsMiwiXFxjQyBcXHRpbWVzIFxcY1Zee1xcb3B9Il0sWzMsMiwiXFxwaXRjaGZvcmsiLDJdLFsxLDIsIlUiXSxbMCwzLCJVIFxcdGltZXMgVV57XFxvcH0iLDJdLFswLDEsIlxccGl0Y2hmb3JrX3svWn0iXV0=
	\[\begin{tikzcd}[ampersand replacement=\&,cramped]
		{Z/\cC \times (\cV/M)^{\op}} \&\& {Z/\cC} \\
		\\
		{\cC \times \cV^{\op}} \&\& \cC
		\arrow["{\pitchfork_{/Z}}", from=1-1, to=1-3]
		\arrow["{U \times U^{\op}}"', from=1-1, to=3-1]
		\arrow["U", from=1-3, to=3-3]
		\arrow["\pitchfork"', from=3-1, to=3-3]
	\end{tikzcd}\cdpunct[8pt]{.}\]
	The functor $U \colon Z/\cC \rightarrow \cC$ creates all limit and lifts the pullback power axiom directly.
\end{proof}

Finally, we specialize to the case of a symmetric closed monoidal category for characterizing all enriched slice model categories.

\begin{theorem}
	Let $(\cV,\otimes, I)$ be a closed symmetric monoidal model category and let $\ucC$ be a $\ucV$-enriched model category with the underlying category $\cC$. Let $(M, \mu, \eta)$ be a commutative monoid object in $\cV$ and let $(Z,\lambda)$ be a left $M$-module in $\cC$ relative to tensoring. Then, it follows that both $Z/\ucC$ and $\ucC/Z$ are $\ucV/M$-enriched model categories.
\end{theorem}


